% ID: SEL-D0047-OA-E02
% ESTADO: validada. S048.
\textbf{a)} $B(0)=-6$ y $B(10)=2000-3600+1620-6=14$. Al inicio hay pérdidas de $6\,000$ euros y al final, un beneficio de $14\,000$ euros.

\textbf{b)} $B'(t)=6(t-3)(t-9)$. Su signo es positivo entre $0$ y $3$, negativo entre $3$ y $9$, y positivo entre $9$ y $10$. Para obtener los extremos absolutos, comparamos ambos puntos críticos y los extremos del intervalo:
\[
\begin{array}{c|rrrr}
t&0&3&9&10\\\hline
B(t)&-6&210&-6&14
\end{array}
\]
El máximo es de $210\,000$ euros en $t=3$. El mínimo es de $-6\,000$ euros y se alcanza en $t=0$ y $t=9$.
